% Full flat fine inverse: fragment for Notebook 09, no standalone preamble.
% Requires amsmath/amssymb/amsthm and the existing \norm, \diag macros.
\section{The full flat source inverse}
\label{inv-section}

We now control the fine inverse itself, rather than only the minimized
coarse action or its retained-coordinate covariance. Fine spacing is
$\eta=M^{-1}$, where $M$ is odd, and the coarse physical spacing is one.
All fine pairings use volume weight $\eta^d$, while coarse pairings use
unit counting measure. Let $D$ be the fine scalar gradient,
$\mathcal L=D^*D$ its positive scalar Laplacian, and $K$ the fine
Maxwell precision. Thus $KD=0$ and $K+DD^*=\mathcal L I$.
Let $Q=Q_M^{\rm st}$ be the normalized one-form straight-string average
and $Q_0$ the centered scalar average, so that $QD=d_cQ_0$.

The scalar source gauge is
\begin{equation}\label{inv-source-projection}
 G'=(\mathcal L+Q_0^*a'Q_0)^{-1},\qquad
 \mathcal P=G'Q_0^*(Q_0G'^2Q_0^*)^{-1}Q_0G',\qquad R=I-\mathcal P.
\end{equation}
The scalar projection $\mathcal P$ must not be confused with a Fourier
denominator $P_r$ below. The positive scalar penalty $a'$ in this
definition is distinct from the one-form penalty $a$ in
\begin{equation}\label{inv-fine-precision}
 \mathcal A_a=K+DRD^*+Q^*aQ,\qquad G_0=\mathcal A_a^{-1}.
\end{equation}
We first allow a fixed positive $a$, uniformly in a fixed interval
$0<a_-\le a\le a_+<\infty$. The source convention of Raw I,
lines 13583--13593, is $a=\ell^{-2}$; in the present homogeneous
unit-scale normalization it is exactly $a=1$.

\begin{theorem}\label{inv-theorem}
The inverse in \eqref{inv-fine-precision} exists on the homogeneous
infinite lattice and every compatible periodic regulator. Its
physical-cell kernel satisfies, for fixed positive $C,\rho$,
\begin{equation}\label{inv-kernel-bound}
 \norm{(\mathcal D_\alpha G_0\mathcal D_\beta)(n,m)}
 _{\mathrm{cell},2\to\mathrm{cell},2}
 \le Ce^{-\rho|n-m|_1},\qquad \alpha+\beta\le2.
\end{equation}
Here $n,m$ are unit coarse cells, each $\mathcal D$ is a specified
composition of fine forward or backward differences of its indicated
order, and fixed curl or divergence combinations are included.
Constants are independent of odd $M$, volume, and $a\in[a_-,a_+]$.
The cell Hilbert norm is the fine physical-volume norm, including the
fixed orientation and color fibers. At every smaller positive exponent,
\eqref{inv-kernel-bound} gives a uniform exponentially weighted
cell-Schur bound. The periodic kernels are exact periodizations and
obey the same bounds in periodic distance.

If $q_M$ denotes the physical Maxwell minimum, then the penalized
coarse inverse and normalized extension satisfy
\begin{equation}\label{inv-coarse-result}
 q_0:=(QG_0Q^*)^{-1}=q_M+aI,
 \qquad G_0Q^*q_0=H_{\rm src},
\end{equation}
where $H_{\rm src}$ is the exact Maxwell minimizing extension in the
matched scalar source gauge $RD^*H_{\rm src}=0$.
\end{theorem}

The theorem is specific to the homogeneous flat source definitions.
It is not a theorem for arbitrary nonflat backgrounds, Dirichlet
domains, or mixed retained histories. Fine derivative bounds in
\eqref{inv-kernel-bound} are cell $L^2$ operator bounds, not
pointwise regularity assertions.

\subsection{Exact aliases and the penalized Feynman inverse}

Let $I_M=\{-(M-1)/2,\ldots,(M-1)/2\}$,
$p_r=z+2\pi r$, and $r\in I_M^d$. Use
\[
 u_{r,\mu}^{\pm}=M(e^{\pm ip_{r,\mu}/M}-1),\quad
 P_r=\sum_\mu u_{r,\mu}^+u_{r,\mu}^-,\quad
 d_\mu^\pm=e^{\pm iz_\mu}-1,\quad b_\mu=d_\mu^+d_\mu^-.
\]
Write $a_r=a_M(p_r)$ for the centered scalar averaging multiplier
and $\omega_{r,\mu}$ for the one-form multiplier, including its
additional longitudinal string factor. Superscript minus means
evaluation at the negative full fine momentum; it agrees with
complex conjugation on real momenta only. String telescoping gives
\begin{equation}\label{inv-cochain}
 \omega_{r,\mu}u_{r,\mu}^+=a_r d_\mu^+,
 \qquad \omega_{r,\mu}^-u_{r,\mu}^-=a_r^-d_\mu^-.
\end{equation}
Alias space has counting Hilbert norm. With physical-volume Parseval
normalization, no additional power of $M$ occurs in an alias covariance
sum.

Work initially in a real nonzero coarse fiber. The scalar source
projection is the orthogonal projection onto
$\operatorname{ran}(\mathcal L^{-1}Q_0^*)$. Therefore
\begin{equation}\label{inv-projection-alias}
 \mathcal P=\mathcal L^{-1}Q_0^*W_2^{-1}Q_0\mathcal L^{-1},
 \qquad W_2=\sum_r\frac{a_ra_r^-}{P_r^2}.
\end{equation}
This agrees with \eqref{inv-source-projection}: multiplying by the
positive scalar penalty changes the coarse parametrization of
$\operatorname{ran}(G'Q_0^*)$, but not that range. Its constant scalar
mode will be treated through the original penalized definition.

Define
\[
 G_F=\diag_r(P_r^{-1}I),\qquad
 W=QG_FQ^*=\diag_\mu(c_\mu),\qquad
 c_\mu=g_\mu/P_0+R_\mu,
\]
\[
 g_\mu=\omega_{0,\mu}\omega_{0,\mu}^-,\qquad
 R_\mu=\sum_{r\ne0}\frac{\omega_{r,\mu}\omega_{r,\mu}^-}{P_r}.
\]
The Feynman operator with the one-form penalty is
$\mathcal B=\mathcal L I+Q^*aQ$. Its inverse obeys
\begin{equation}\label{inv-F-woodbury}
 F:=\mathcal B^{-1}
 =G_F-G_FQ^*(a^{-1}I+W)^{-1}QG_F.
\end{equation}
Put
\begin{equation}\label{inv-Da}
 \widehat R_\mu=R_\mu+a^{-1},\qquad
 \mathcal D_\mu=g_\mu+P_0\widehat R_\mu,\qquad
 k_\mu=\mathcal D_\mu^{-1}.
\end{equation}
Then $(a^{-1}I+W)^{-1}=P_0\diag(k_\mu)$ and
\begin{equation}\label{inv-F-alias}
 F_{r\mu,s\nu}=\delta_{\mu\nu}
 \left[\frac{\delta_{rs}}{P_r}
 -\frac{P_0\omega_{r,\mu}^-\omega_{s,\mu}}
              {P_rP_s\mathcal D_\mu}\right].
\end{equation}
The central blocks cancel exactly to
\begin{align}
 F_{0\mu,0\nu}
 &=\delta_{\mu\nu}\widehat R_\mu/\mathcal D_\mu,
 \label{inv-F-central}\\
 F_{0\mu,s\nu}
 &=-\delta_{\mu\nu}
   \frac{\omega_{0,\mu}^-\omega_{s,\mu}}{P_s\mathcal D_\mu}
       &&(s\ne0),\nonumber\\
 F_{r\mu,0\nu}
 &=-\delta_{\mu\nu}
   \frac{\omega_{r,\mu}^-\omega_{0,\mu}}{P_r\mathcal D_\mu}
       &&(r\ne0).\nonumber
\end{align}
Use \eqref{inv-F-alias} only with both aliases high, and
\eqref{inv-F-central} otherwise. These expressions contain no
central massless denominator.

\subsection{The scalar projection and the second Woodbury identity}

The Hodge identity and $R=I-\mathcal P$ give
\begin{equation}\label{inv-rank-subtraction}
 \mathcal A_a=\mathcal B-vW_2^{-1}v^-,\qquad
 v_{r,\mu}=u_{r,\mu}^+a_r^-/P_r,
 \quad v_{r,\mu}^-=a_ru_{r,\mu}^-/P_r.
\end{equation}
Here $v^-$ is a row, the analytic transpose at negative arguments.
On real fibers it is the Hilbert adjoint of $v$.

The real fine precision is positive definite. Indeed, a vector in its
kernel has zero curl, zero coarse average, and zero source-gauge term.
The zero-curl field is $D\lambda+c$, with $c$ a constant harmonic
one-form. Taking the coarse mean of $Q(D\lambda+c)=0$ gives $c=0$;
subtracting a scalar constant makes $Q_0\lambda=0$. Then
$\mathcal L\lambda\in\operatorname{ran}R$, while the source-gauge
condition says $R\mathcal L\lambda=0$. Thus $D\lambda=0$.
This proves finite-regulator positivity with the harmonic sector
retained. The subtractive Woodbury identity therefore gives
\begin{equation}\label{inv-G-woodbury}
 G_0=F+Fv\,(W_2-v^-Fv)^{-1}v^-F.
\end{equation}

Equation \eqref{inv-cochain} implies the three exact contractions
\begin{equation}\label{inv-three-contractions}
 v^-G_Fv=W_2,\qquad QG_Fv=W_2d^+,
 \qquad v^-G_FQ^*=W_2d^-.
\end{equation}
For example, the $\mu$ component of the middle term is
$\sum_r\omega_{r,\mu}u_{r,\mu}^+a_r^-/P_r^2=W_2d_\mu^+$.
Let $\mathfrak a=a_0a_0^-$. The central cochain identity is
$\sum_\mu b_\mu/g_\mu=P_0/\mathfrak a$. Define
\begin{equation}\label{inv-Ea}
 \mathcal E_a=\mathfrak a^{-1}
 -\sum_\mu\frac{b_\mu\widehat R_\mu}{g_\mu\mathcal D_\mu}.
\end{equation}
Subtracting $1/\mathcal D_\mu-1/g_\mu$ proves
\begin{equation}\label{inv-Ea-identity}
 \sum_\mu b_\mu k_\mu=P_0\mathcal E_a.
\end{equation}
It follows that
\begin{align}
 W_2-v^-Fv
 &=W_2^2d^-(a^{-1}I+W)^{-1}d^+
  =W_2^2P_0^2\mathcal E_a.
 \label{inv-woodbury-denominator}
\end{align}

Introduce the scalar second-order remainder and its canceled denominator
\begin{equation}\label{inv-Ds2}
 R_{s,2}=\sum_{r\ne0}\frac{a_ra_r^-}{P_r^2},\qquad
 D_{s,2}=\mathfrak a+P_0^2R_{s,2},\qquad
 W_2=D_{s,2}/P_0^2.
\end{equation}
Thus \eqref{inv-woodbury-denominator} is
$D_{s,2}^2\mathcal E_a/P_0^2$. In particular it contains $W_2^2$,
not $W_2$.

\subsection{The final canceled inverse}

Set $\beta=(P_0/D_{s,2})Fv$. For a high alias,
\begin{equation}\label{inv-beta-high}
 \beta_{r,\mu}
 =\frac{P_0u_{r,\mu}^+a_r^-}{P_r^2D_{s,2}}
 -\frac{\omega_{r,\mu}^-k_\mu d_\mu^+}{P_r},
 \qquad r\ne0.
\end{equation}
For the central entry use
\[
 \omega_{0,\mu}^-k_\mu d_\mu^+
 =(u_{0,\mu}^+/a_0)g_\mu k_\mu,
\]
\[
 \mathfrak a/D_{s,2}=1-P_0^2R_{s,2}/D_{s,2},\qquad
 g_\mu k_\mu=1-P_0\widehat R_\mu k_\mu.
\]
The resulting exact central formula is
\begin{equation}\label{inv-beta-central}
 \beta_{0,\mu}=
 \frac{u_{0,\mu}^+}{a_0}
 \left(\widehat R_\mu k_\mu
                  -\frac{P_0R_{s,2}}{D_{s,2}}\right).
\end{equation}
The row $\beta^-=(P_0/D_{s,2})v^-F$ is its analytic transpose:
in the high formula it has entries
$P_0a_ru_{r,\mu}^-/(P_r^2D_{s,2})
-d_\mu^-k_\mu\omega_{r,\mu}/P_r$; centrally use
$u_{0,\mu}^-/a_0^-$ times the same bracket.
Equations \eqref{inv-G-woodbury} and
\eqref{inv-woodbury-denominator} now become
\begin{equation}\label{inv-final-formula}
 \boxed{\displaystyle G_0=F+\beta\mathcal E_a^{-1}\beta^-.}
\end{equation}
Every central pole has canceled. The update has a positive sign, and
the second term in the bracket of \eqref{inv-beta-central} has a
negative sign. The formula removes the entire complex cone $P_0=0$.

\subsection{A strip and constants uniform in refinement}

Choose $\pi<K_0<2\pi$ and set
\[
 c_* =\frac{\sin(K_0/2)}{K_0/2},\qquad
 H_* =\frac{4(2\pi-K_0)^2}{\pi^2}.
\]
On the enlarged real cube $|x_\mu|\le K_0$, the alias estimates of
Notebook 07 give
\[
 \mathfrak a\ge c_*^{2d},\quad
 g_\mu\ge g_*:=c_*^{2d+2},\quad
 0\le R_\mu\le H_*^{-1},\quad
 0\le R_{s,2}\le H_*^{-2},\quad P_0\le dK_0^2,
\]
with $g_\mu,\mathfrak a\le1$. Consequently
\[
 g_*\le\mathcal D_\mu\le
 D_*:=1+dK_0^2(H_*^{-1}+a_-^{-1}),\qquad
 c_*^{2d}\le D_{s,2}\le1+d^2K_0^4H_*^{-2}.
\]
For $x\ne0$, \eqref{inv-Ea-identity} gives
\begin{equation}\label{inv-Ea-lower}
 \mathcal E_a(x)=
 \frac{\sum_\mu b_\mu(x)/\mathcal D_\mu(x)}{P_0(x)}
 \ge c_*^2/D_*.
\end{equation}
Here $\sum b_\mu\ge c_*^2|x|^2$ and $P_0\le|x|^2$.
At zero the defining expression gives $\mathcal E_a(0)=1$.

On an initial fixed complex neighborhood, high denominators obey
$|P_r(z)|\ge c|r|^2$ for $r\ne0$, including the edge aliases.
Finite alias Parseval bounds the sums of absolute paired scalar and
one-form averaging factors independently of $M$. Hence $R_\mu$,
$R_{s,2}$, and the central multipliers are uniformly holomorphic and
bounded. The added $a^{-1}$ is bounded by the hypothesis on $a$.

Use nested real cubes with positive margins. Uniform Cauchy estimates
first allow an imaginary shrinking that keeps $\mathfrak a,g_\mu,
\mathcal D_\mu,D_{s,2}$ separated from zero. The expression
\eqref{inv-Ea} is then holomorphic and uniformly bounded. A second
shrinking, using \eqref{inv-Ea-lower}, keeps $\mathcal E_a$ separated
from zero. This gives one positive strip width independent of $M,a$,
with the real fundamental cube in the interior of each enlarged chart.
The inverse of $a_0$ is also bounded because $a_0a_0^-=\mathfrak a$
is separated from zero and both factors are bounded. Real positivity
has thus been combined with uniform analytic estimates; it is not
being used alone to claim complex invertibility.

\subsection{Zero mode and the independent $M=1$ check}

At coarse momentum zero, $a_0=\omega_{0,\mu}=1$, all high
$a_r,\omega_{r,\mu}$ vanish, and
\[
 R_\mu=R_{s,2}=0,\qquad
 \mathfrak a=g_\mu=\mathcal D_\mu=D_{s,2}=\mathcal E_a=1,
 \qquad\beta=0.
\]
The canceled formula gives
\begin{equation}\label{inv-zero-mode}
 (G_0)_{0\mu,0\nu}=a^{-1}\delta_{\mu\nu},\qquad
 (G_0)_{r\mu,s\nu}=\delta_{rs}\delta_{\mu\nu}/P_r
       \quad(r,s\ne0),
\end{equation}
with zero mixed blocks. The original source projection is zero on the
constant scalar mode and identity on high scalar aliases. Hence the
central precision is $aI$, and the high precision is $P_rI$, exactly
as in \eqref{inv-zero-mode}. No scalar constant Laplacian was inverted.

For $M=1$, $Q=I$, $R=0$, and
\[
 \mathcal D=1+P_0/a,\quad k=\mathcal E_a=a/(P_0+a),\quad
 F=(P_0+a)^{-1}I,\quad\beta=d^+/(P_0+a).
\]
Thus
\begin{equation}\label{inv-M1}
 G_0=(P_0+a)^{-1}I+\frac{d^+d^-}{a(P_0+a)},
\end{equation}
the direct inverse of $(P_0+a)I-d^+d^-$. This checks the sign,
penalty normalization, and rank-one coefficient. The value $a^{-1}I$
at zero also shows why a uniform lower bound for $a$ is essential.

\subsection{Alias bounds with two fine derivatives}

On the common strip,
\[
 |P_r|\ge c|r|^2\ (r\ne0),\quad
 |u_{r,\mu}^{\pm}|\le C(1+|r|),\quad
 \sum_r|a_r^\pm|^2+
       \sum_{r,\mu}|\omega_{r,\mu}^\pm|^2\le C_d.
\]
The high-high block of $F$ is a bounded diagonal operator plus one
rank-one term for each orientation. Its high rank vectors are
$\omega_{r,\mu}^{\pm}/P_r$, and its mixed row and column use the
same vectors. Their alias $L^2$ norms are bounded. The central block
is bounded separately. Moreover \eqref{inv-beta-high} gives
\begin{equation}\label{inv-beta-estimate}
 |\beta_{r,\mu}|\le
 C\left(\frac{|a_r^-|}{|r|^3}
             +\frac{|\omega_{r,\mu}^-|}{|r|^2}\right),\qquad r\ne0.
\end{equation}
The central beta is bounded by \eqref{inv-beta-central}.
Therefore $F$ and $G_0$ have uniformly bounded alias Hilbert operator
norms. The norm estimate for an analytic rank-one matrix $v w^T$
uses $\norm{v}_2\norm{w}_2$, without conjugating its analytic parameters.

With $\alpha$ differences on the output and $\beta$ on the input,
$\alpha+\beta\le2$, the high diagonal of $F$ grows at most as
$|r|^{\alpha+\beta-2}$, which is bounded. Each rank vector has at
most two extra powers of $|r|$ multiplying
$\omega_{r,\mu}^{\pm}/P_r$, so its $L^2$ norm is still bounded by
Parseval. The mixed blocks have the same property. For the update in
\eqref{inv-final-formula}, two derivatives on a beta column are bounded,
in the worst case, by $C(|a_r^-|/|r|+|\omega_{r,\mu}^-|)$.
This is square summable uniformly, as is the corresponding row.
Every derivative sandwich of total order at most two is consequently
holomorphic and uniformly bounded in alias operator norm. The number
of derivative-orientation components is fixed in $d$.

\subsection{Reconstruction, periodic gluing, and the actual inverse}

For within-cell sites $s\in I_M^d$ use
$\mathcal F_M(z)_{sr}=e^{i(z+2\pi r)\cdot s/M}$.
For real $z$ it is unitary from alias counting $L^2$ into the cell
norm $M^{-d}\sum_s|v_s|^2$. For complex $z=x+iy$, it differs from
that unitary by $\diag_s(e^{-y\cdot s/M})$. Its norm and the norm of
its holomorphic inverse are thus uniformly bounded. Tensoring with
the orientation identity, the reconstructed symbol is
\[
 (G_0)_{\rm cell}(z)=\mathcal F_M(z)G_0(z)\mathcal F_M(z)^{-1}.
\]
All derivative sandwiches inherit its uniform strip bounds.

The full real inverse, rather than each separated central term, is
periodic up to alias permutation. Reconstruction cancels the
permutation; wrapping an alias gives $e^{i2\pi s_\nu}=1$. The
reconstructed symbol is therefore periodic. The enlarged-chart
holomorphic formulas agree on real overlaps and glue by the identity
theorem. Fourier contour shifting proves \eqref{inv-kernel-bound}.
Summing at a smaller exponent proves the weighted Schur bounds.
These are operator-valued contour estimates, so no factor growing
with the within-cell dimension is introduced.

For each fixed $M$, fiberwise multiplication, with the zero sector
as in \eqref{inv-zero-mode}, verifies that this bounded kernel is
the inverse of \eqref{inv-fine-precision} on the infinite homogeneous
Hilbert space. The real inverse bound yields a uniform positive fine
spectral gap. Sampling the same symbol on a compatible periodic
momentum grid gives exactly the finite source inverse. Absolute
summability identifies its kernel with the periodization; periodic
distance bounds its weighted sums by the infinite sums. This includes
short periods with their proper oriented-cell multiplicities.

Finally, at fixed coarse datum $B$, the Maxwell energy is at least
$\langle B,q_MB\rangle$, the source gauge term is nonnegative, and
the penalty is the constant $a\norm{B}^2$. The source-gauge minimizer
$H_{\rm src}B$ attains both minima. Therefore
\[
 \inf_{QA=B}\langle A,\mathcal A_aA\rangle
   =\langle B,(q_M+aI)B\rangle.
\]
The constrained Gaussian inverse formula proves
$q_0=q_M+aI$ and, by strict positivity and uniqueness,
$G_0Q^*q_0=H_{\rm src}$. In the source convention
$QH_0=\ell^{-1}$, the normalized extension is $H_0\ell$.
This proves \eqref{inv-coarse-result} and completes the theorem.

\begin{remark}\label{inv-scope}
The homogeneous source passage fixes $a=\ell^{-2}$ and hence $a=1$
here. Raw I, lines 23879--23881, separately introduces an $a_n$
sequence and refers to CMP99b (3.23)--(3.24) for an unprinted recursion.
The theorem does not guess that sequence: an application using it
must verify the fixed positive interval after its measure and scale
conversion. The scalar penalty $a'$ in the projection remains a
different object. General nonflat and retained-history identifications
are not consequences of the homogeneous Fourier proof.
\end{remark}
